% from ega4-8-fr.tex
\oldpage[IV-3]{19}
\begin{corollary}[8.4.3]
\label{IV.8.4.3-fr}
Supposons vérifiée l'une des conditions \emph{a)}, \emph{b)} de (8.4.2),
l'espace sous-jacent à $S$ étant noethérien~; alors, pour que $S$ soit
irréductible, il faut et il suffit qu'il existe $\lambda$ tel que $S_\mu$ le
soit pour tout $\mu\geq\lambda$.
\end{corollary}

\begin{proposition}[8.4.4]
\label{IV.8.4.4-fr}
Supposons qu'il existe un $\alpha$ tel que $S_\alpha$ soit quasi-compact et
que, pour $\lambda\leq\mu$, $u_{\lambda\mu}:S_\mu\to S_\lambda$ soit dominant.
Alors, pour que $S$ soit connexe, il faut et il suffit qu'il existe $\lambda$
tel que $S_\mu$ le soit pour tout $\mu\geq\lambda$.
\end{proposition}

\begin{proof}
La condition est suffisante en vertu de (8.4.1)~; d'autre part, on a vu
(8.3.8, (i)) que $u_\lambda:S\to S_\lambda$ est dominant pour $\lambda$ assez
grand, donc, si $S$ est connexe, il en est de même de $S_\lambda$, puisque
$u_\lambda(S)$ est dense dans $S_\lambda$ et connexe.
\end{proof}

\begin{corollary}[8.4.5]
\label{IV.8.4.5-fr}
Soient $k$ un corps, $X$ un $k$-préschéma quasi-compact. Pour que $X$ soit
géométriquement connexe (4.5.2), il faut et il suffit que, pour toute extension
finie et séparable $K$ de $k$, $X\otimes_k K$ soit connexe.
\end{corollary}

\begin{proof}
La condition est trivialement nécessaire. Pour voir qu'elle est suffisante,
nous devons prouver que si $\Omega$ est une clôture algébrique de $k$,
$X\otimes_k\Omega$ est connexe (4.5.1). Or, $\Omega$ est limite inductive
filtrante des sous-extensions finies $k'$ de $k$, et pour
$k\subset k'\subset k''\subset\Omega$, le morphisme
$X\otimes_k k''\to X\otimes_k k'$ est surjectif. On est donc ramené, en vertu
de (8.4.4), à prouver que $X\otimes_k k'$ est connexe pour toute extension
finie $k'$ de $k$. Mais si $K$ est la plus grande extension séparable contenue
dans $k'$, le morphisme $X\otimes_k k'\to X\otimes_k K$ est fini, surjectif et
radiciel, donc (2.4.5) un homéomorphisme, et comme $X\otimes_k K$ est connexe
par hypothèse, il en est de même de $X\otimes_k k'$.
\end{proof}

\begin{remark}[8.4.6]
\label{IV.8.4.6-fr}
\begin{enumerate}[label=(\roman*)]
  \item La démonstration de (8.4.2) montre que la conclusion de cette
  proposition est valable si l'on suppose que l'espace sous-jacent à $S$ est
  noethérien, qu'il existe $\alpha$ tel que $S_\alpha$ soit quasi-compact, et
  enfin que les $u_\lambda:S\to S_\lambda$ sont dominants.

  \item Par contre, la conclusion de (8.4.2) peut être en défaut lorsque les
  $u_\lambda$ ne sont pas dominants pour $\lambda$ assez grand, même lorsque
  les $S_\lambda$ et $S$ sont noethériens, ainsi que le montre l'exemple
  suivant. Prenons pour ensemble d'indices $\mathbf N$, tous les $S_n$ étant
  égaux à
  $\operatorname{Spec}(A\times K)=\operatorname{Spec}(A)\amalg
  \operatorname{Spec}(K)$, où $K$ est un corps, $A$ une $K$-algèbre
  quelconque, et tous les morphismes $u_{n,n+1}$ étant égaux au même morphisme
  correspondant à l'homomorphisme $(x,y)\mapsto(j(y),y)$ de $A\times K$ dans
  lui-même, où $j:K\to A$ est l'homomorphisme canonique. On vérifie aisément
  que la limite inductive de ce système d'anneaux est $K$, l'homomorphisme
  canonique $u_n$ correspondant à la seconde projection $A\times K\to K$. On
  voit donc que $S=\operatorname{Spec}(K)$ est irréductible bien qu'aucun des
  $S_n$ ne soit connexe.
\end{enumerate}
\end{remark}

\subsection{Modules de présentation finie sur une limite projective de préschémas}
\label{subsection:IV.8.5-fr}

\begin{env}[8.5.1]
\label{IV.8.5.1-fr}
Nous conservons toujours les notations de (8.2.2)~; nous nous bornerons en
outre au cas où $S_0$ est l'un des $S_\lambda$, auquel on peut toujours se
ramener.

\emph{Lorsque, dans cette section, nous considérerons une famille
$(\mathcal F_\lambda)$, où, pour tout $\lambda\in L$, $\mathcal F_\lambda$ est
un $\mathcal O_{S_\lambda}$-Module, il sera sous-entendu que cette famille
vérifie la condition}
\begin{equation}
  \mathcal F_\mu=u_{\lambda\mu}^*(\mathcal F_\lambda)
  \qquad\text{pour }\lambda\leq\mu.
  \tag{8.5.1.1}\label{IV.8.5.1.1-fr}
\end{equation}

